%% =====================================================================
%%  SUPPLEMENTARY MATERIAL
%%  "Minimising station visits in automated pick-and-pass warehouses:
%%   a correlated storage assignment model validated at industrial scale"
%%  International Journal of Systems Science: Operations & Logistics
%%
%%  This file holds the statistical arguments and secondary results the
%%  main article references (S1 calibration, S2 sensitivity, S3 paired
%%  tests, S4 implementation record, S5 proofs, S6 Company A instance,
%%  S7 industrial beta grid, S8 convergence record). It compiles
%%  standalone: it needs only IJSSOL_CSLAP_v1.bib alongside it. Cross-references to the main article are written out
%%  in full, since the two documents compile separately.
%% =====================================================================
\PassOptionsToPackage{table}{xcolor}
\documentclass[11pt]{article}

\usepackage[a4paper,margin=2.5cm]{geometry}
\usepackage{amsmath,amssymb,amsfonts}
\usepackage{mathtools}
\usepackage{booktabs}
\usepackage{float}
\usepackage{xcolor}
\usepackage{tabularx}
\usepackage{multirow}
\usepackage{url}
\usepackage{hyperref}

\usepackage[natbibapa,nodoi]{apacite}
\setlength\bibhang{12pt}

% Number sections, tables and figures with an S prefix, as T&F expect
% for supplementary items.
\renewcommand{\thesection}{S\arabic{section}}
\renewcommand{\thetable}{S\arabic{table}}
\renewcommand{\thefigure}{S\arabic{figure}}

% The main article uses \tbl{} and \tabnote{} from interact.cls; supply
% equivalents so the table bodies below are unchanged from the article.
\newcommand{\tbl}[1]{\caption{#1}}
\newcommand{\tabnote}[1]{\vspace{2pt}\footnotesize\raggedright #1\par}

\title{Supplementary material:\\
Calibration, proofs and implementation record}
\author{Ermal Belul \and Marwane Bouznif \and Dritan Nace \and Antoine Jouglet}
\date{}

\begin{document}
\maketitle

\noindent
This supplement holds the statistical arguments and the secondary results behind the main
article; every definition, notation and method statement remains in the main text.
Sections~\ref{sup:capsweep} to~\ref{sup:tests} concern the greedy clustering heuristic of
Section~4.1 of the main article and the synthetic benchmark of its Section~5: whether the
community size bound $\beta$ scales with catalogue size ($N$) or with station slot capacity
($\zeta$), how sensitive the solution is to each threshold ($\phi, \gamma, r, \beta$), and the
paired statistical tests behind the benchmark. Section~\ref{sup:impl} records the
implementation, hardware and parameter settings of every method; Section~\ref{sup:proofs}
proves the two bounding claims of the main article; Section~\ref{sup:instance} records the
Company~A instance; Section~\ref{sup:betasearch} reports the full industrial $\beta$ search; and
Section~\ref{sup:cgrecord} describes the released per-instance column-generation convergence
record.

\section{Community bound scaling across warehouse geometries}\label{sup:capsweep}
The three filters of Section~4.1 of the main article are indexed on the catalogue and the
community bound on the station. That asymmetry needs evidence, because the benchmark of
Section~5 cannot supply it: its generator ties the station count to the catalogue, so
$\zeta=100$ at every size from 500 SKUs upward and the ratio $\beta/\zeta$ is constant across
the whole family. Catalogue size and slot capacity are confounded there, and no result measured
on it can separate them.

We therefore built a second family in which the two are not confounded. It uses the same
generator as Section~5 of the main article (the common-itemset procedure of
\citet{zhang2019cie} at correlation degree $0.7$, with order size drawn uniformly from the integers $5$ to $15$ and
proportional order volumes) and changes one thing: the number of stations is set independently
of the catalogue rather than as $\max(5,\lfloor N/100\rfloor)$. Two catalogue sizes,
$N\in\{500,1000\}$, are crossed with four station counts chosen so that the same slot
capacities $\zeta\in\{20,25,50,100\}$ occur at both sizes. Ten instances per cell give 80 in
all, generated on seeds $2001$ to $2010$, disjoint from the seeds of the benchmark. Station
labels are drawn from a separate random stream, so at a fixed seed the order structure is
identical across the four geometries and only the station configuration changes. Any difference
between cells is therefore attributable to geometry alone.

Each instance is solved at every $\beta$ on the grid
$\{2,3,5,8,10,12,15,18,20,25,30,40,50,75,100\}$ that falls at or below its own $\zeta$, which
gives 980 runs. All 980 returned a feasible layout. Table~\ref{tab:capsweep} reports the
minimising bound in each cell. Both instance families, with the generator settings and seeds
needed to reconstruct them exactly, accompany the article.

\begin{table}[H]
\centering
\tbl{Visit-minimising community bound $\beta^\ast$ across warehouse geometries, on 80 instances disjoint from the benchmark. The same $\zeta$ appears at both catalogue sizes.}
{\small\begin{tabular*}{\textwidth}{@{\extracolsep{\fill}}rrrrrr@{}}
\toprule
$\zeta$ & $N$ & $|S|$ & $\beta^\ast$ & $\beta^\ast/\zeta$ & $\Delta$ visits vs.\ $\beta=15$ \\
\midrule
20  & 500   & 25 & 12 & 0.60 & $-0.72$\% \\
20  & 1{,}000 & 50 & 12 & 0.60 & $-0.45$\% \\
25  & 500   & 20 & 15 & 0.60 & $0.00$\% \\
25  & 1{,}000 & 40 & 15 & 0.60 & $0.00$\% \\
50  & 500   & 10 & 30 & 0.60 & $-0.87$\% \\
50  & 1{,}000 & 20 & 30 & 0.60 & $-0.84$\%\textsuperscript{$\ast$} \\
100 & 500   & 5  & 40 & 0.40 & $-1.01$\%\textsuperscript{$\ast$} \\
100 & 1{,}000 & 10 & 40 & 0.40 & $-1.06$\%\textsuperscript{$\ast$} \\
\bottomrule
\end{tabular*}}
\tabnote{$\beta^\ast$ minimises mean visits in its cell over the swept values, which include 15. The last column is the change from 15 to $\beta^\ast$, so a negative entry means $\beta^\ast$ uses fewer visits and a zero entry means 15 is itself the minimiser, as it is in both $\zeta=25$ cells. \textsuperscript{$\ast$}Separates from 15 at $p<0.05$ under a Holm-corrected Wilcoxon signed-rank test over the ten instances of the cell. The unstarred non-zero cells fail that test on within-cell variance rather than on effect size: at $\zeta=50$, $N=500$ the point effect of $-0.87$\% exceeds that of the significant cell beneath it.}
\label{tab:capsweep}
\end{table}

At a fixed slot capacity $\beta^\ast$ is the same at both catalogue sizes in all four cells,
whereas at a fixed catalogue size it moves with $\zeta$ across the whole range. Over the eight
cells the Spearman rank correlation between $\beta^\ast$ and $\zeta$ is $1.000$: the two
orderings agree exactly. The Pearson correlation on the raw values is $0.962$, the shortfall
measuring only that the dependence is not linear, as the next paragraph shows. Against $N$
both measures return $0.000$. The bound therefore scales with the slot capacity of a station
and is insensitive to the length of the catalogue.

The dependence is monotone but not proportional. The minimising ratio $\beta^\ast/\zeta$ is
$0.60$ at $\zeta\le 50$ and falls to $0.40$ at $\zeta=100$, so a single coefficient cannot be
optimal everywhere. We adopt $0.4$ for two reasons. It is the minimising value at $\zeta=100$,
the capacity of every benchmark instance of 500 SKUs and above and of the sizes at which the
method is intended to be used. Where it is not the minimiser it errs low, which matters because
larger communities are harder to place within the workload constraint, as Section~6 of the main
article shows on a site whose stations are not interchangeable. The price of that choice is the
excess of the layout built at $\beta=\lfloor 0.4\zeta+\tfrac12\rfloor$ over the layout built at
the cell's own $\beta^\ast$, paired instance by instance and averaged over the ten instances of
the cell. It is nil at $\zeta=100$, where $0.4\zeta$ is itself the minimiser, and runs from
$0.15$\% at $\zeta=50$, $N=500$ to $1.03$\% at $\zeta=25$, $N=500$ across the six cells below
it. That quantity is not the last column of Table~\ref{tab:capsweep}, which measures $\beta^\ast$
against the fixed value 15 rather than against the rule.

\section{Threshold sensitivity of the clustering heuristic}\label{sup:sensitivity}
Table~\ref{tab:sensitivity} rescales each of the four thresholds of Section~4.1 of the main
article from half to twice its nominal value, one at a time and then all four together, on the
twelve 50-SKU and ten 500-SKU instances. The 1,000- and 2,000-SKU instances are omitted because
four and three instances are too few for an informative interval, although the ordering of the
effects there agrees with what follows. A multiplier applies to the effective nominal value on
the instance, that is, after the threshold's own clamp has been taken (the pre-clamp value comes
from $N$ for the three filters and from $\zeta$ for the community bound), and the scaled result
is clamped again, at 1 for the two frequency floors and at 2 for the community bound. At 50 SKUs
both frequency floors already sit on their clamp of 2, so $\times0.75$ leaves them there and
$\times0.5$ only lowers them to 1; the community bound sits on its clamp of 5, put there by the
floor rather than by $0.4\zeta$, and its downward settings do move it, to 4 and to 2.
Entries are paired mean changes against the nominal setting on the same instance. The
procedure is deterministic: repeating the nominal setting under five hash seeds returns an
identical layout on every instance, so any difference is a real threshold effect and not
run-to-run variation. Reporting follows the conventions for computational experiments with
heuristics \citep{barr1995designing, hooker1995testing}.

The two frequency floors are inert, as rescaling either leaves the visit count unchanged and
the surviving edge set bit-identical at every size, because neither removes a product or pair
that the kept-edge ratio does not already remove. The ratio changes the edge set substantially
but the objective very little, separating from the nominal setting only at $\times2$. The
community bound is the one threshold whose response is not monotone. Halving it costs $0.63\%$
at 500 SKUs and doubling it costs $1.19\%$, while $\times0.75$ and $\times1.5$ are
indistinguishable from the nominal setting, while at 1,000 and 2,000 SKUs the pattern is the
same, with doubling costing $1.84\%$ and $2.27\%$. The exception is $\times1.5$ at 500 SKUs,
which recovers $0.45\%$, consistent with Section~\ref{sup:capsweep} above, where the minimising
ratio at $\zeta=100$ is $0.40$ but rises to $0.60$ at smaller capacities, so a bound half again
as large is not yet past the minimum on every instance. A default whose response rises on both
sides at the two larger sizes, and is flat rather than falling at the smaller one, is indexed on
the right quantity, whereas a default with a monotone response is not. Leaving all four at their
defaults is therefore supported here, not assumed.

\begin{table}[H]
\centering
\tbl{Sensitivity of the clustering heuristic to its four thresholds, on the 22 instances of 50 and 500 SKUs.}
{\small\setlength{\tabcolsep}{4pt}\begin{tabularx}{\textwidth}{@{}X c c@{}}
\toprule
Threshold setting & $\Delta$ visits, 50 SKUs & $\Delta$ visits, 500 SKUs \\
\midrule
Frequency floor $\phi$, $\times0.5$ to $\times2$\textsuperscript{a}     & $0.00$\% & $0.00$\% \\
Co-occurrence floor $\gamma$, $\times0.5$ to $\times2$\textsuperscript{a} & $0.00$\% & $0.00$\% \\
\midrule
Kept-edge ratio $r$, $\times0.5$   & $-0.77$\% [$-2.13$, $+0.58$] & $+0.16$\% [$-0.09$, $+0.41$] \\
Kept-edge ratio $r$, $\times0.75$  & $-0.66$\% [$-1.98$, $+0.65$] & $0.00$\% \\
Kept-edge ratio $r$, $\times1.5$   & $+0.23$\% [$-0.72$, $+1.17$] & $+0.08$\% [$-0.27$, $+0.43$] \\
Kept-edge ratio $r$, $\times2$     & $-0.74$\% [$-2.26$, $+0.78$] & $+1.08$\% [$+0.18$, $+1.98$] \\
\midrule
Community bound $\beta$, $\times0.5$  & $+1.77$\% [$-0.17$, $+3.72$] & $+0.63$\% [$-0.07$, $+1.32$] \\
Community bound $\beta$, $\times0.75$ & $+0.25$\% [$-1.19$, $+1.69$] & $+0.03$\% [$-0.46$, $+0.51$] \\
Community bound $\beta$, $\times1.5$  & $-0.25$\% [$-2.29$, $+1.79$] & $-0.45$\% [$-1.20$, $+0.30$] \\
Community bound $\beta$, $\times2$    & $-0.25$\% [$-1.60$, $+1.10$] & $+1.19$\% [$+0.25$, $+2.13$] \\
\midrule
All four, $\times0.5$ & $+1.15$\% [$-1.14$, $+3.43$] & $+0.59$\% [$-0.08$, $+1.27$] \\
All four, $\times2$   & $-1.10$\% [$-2.57$, $+0.38$] & $+1.78$\% [$+1.02$, $+2.54$] \\
\bottomrule
\end{tabularx}}
\tabnote{Entries are mean paired changes in station visits against the nominal setting, with 95\% confidence intervals; a positive value means more visits. Every setting returned a feasible layout on every instance. \textsuperscript{a}All four settings give bit-identical output, so the rows are collapsed into one.}
\label{tab:sensitivity}
\end{table}

\section{Paired statistical testing of the synthetic benchmark}\label{sup:tests}
Section~5 of the main article compares the methods on mean station visits, gaps, running time
and workload feasibility across the 29 synthetic instances. This section reports the paired
statistical testing behind those comparisons, which is kept out of the main text because its
resolution is bounded by the instance counts and adds little to the readings the main tables
already support.

Instances differ by two orders of magnitude in difficulty, so comparing two methods on their
averages would confound the between-method difference with the between-instance variance. We
therefore compare them instance by instance. For each pair of methods we take the per-instance
difference in station visits and test whether those differences are systematically signed. The
test is the two-sided exact Wilcoxon signed-rank test \citep{wilcoxon1945}, the standard choice
for comparing algorithms across a set of benchmark instances because it imposes no
distributional assumption on the differences \citep{demsar2006}, and ties are discarded. The
reported $p$ is the probability, under the null hypothesis that the paired differences are
symmetric about zero, of observing a signed-rank statistic at least as extreme as the one
obtained, in either direction. Table~\ref{tab:tests} reports the outcome against the
set-variable formulation, the per-instance reference of the main article.

\begin{table}[H]
\centering
\tbl{Paired Wilcoxon signed-rank tests on the 29 synthetic instances, against the set-variable formulation.}
{\small\begin{tabular*}{\textwidth}{@{\extracolsep{\fill}}llrrrr@{}}
\toprule
Comparison & Size (SKUs) & $n$ & Wins & $p$ & Min.\ $p$ \\
\midrule
Binary MILP vs.\ reference & 50      & 12 & 4  & 0.74   & 0.0005 \\
                           & 500     & 10 & 8  & 0.049  & 0.0020 \\
                           & 1{,}000 & 4  & 1  & 0.38   & 0.125  \\
                           & 2{,}000 & 3  & 0  & 0.25   & 0.25   \\
                           & pooled  & 29 & 13 & 0.91   & --     \\
\midrule
CG vs.\ reference          & 50      & 12 & 1  & 0.0010 & 0.0005 \\
                           & 500     & 10 & 2  & 0.0273 & 0.0020 \\
                           & 1{,}000 & 4  & 2  & 1.00   & 0.125  \\
                           & 2{,}000 & 3  & 3  & 0.25   & 0.25   \\
                           & pooled  & 29 & 8  & 0.18   & --     \\
\midrule
GA and SA-C vs.\ reference & 50      & 12 & 0  & 0.0005 & 0.0005 \\
                           & 500     & 10 & 0  & 0.0020 & 0.0020 \\
                           & 1{,}000 & 4  & 0  & 0.125  & 0.125  \\
                           & 2{,}000 & 3  & 0  & 0.25   & 0.25   \\
\midrule
Heuristic vs.\ reference   & 50      & 12 & 0  & 0.0005 & 0.0005 \\
                           & 500     & 10 & 0  & 0.0020 & 0.0020 \\
                           & 1{,}000 & 4  & 0  & 0.125  & 0.125  \\
                           & 2{,}000 & 3  & \textbf{3}  & 0.25   & 0.25   \\
\bottomrule
\end{tabular*}}
\tabnote{``Wins'' counts instances on which the first-named method returns fewer visits. ``Min.\ $p$'' is the smallest two-sided $p$ attainable at that sample size, $2^{1-n}$, and where $p$ equals it the test is saturated rather than failed, so the win count carries the evidence. The GA, SA-C and Heuristic rows coincide at 50, 500 and 1{,}000 SKUs because none of the three wins a single paired instance there. The Heuristic separates from them only at 2{,}000 SKUs, where it wins all three.}
\label{tab:tests}
\end{table}

No test at the two larger sizes can resolve anything. The resolution of a signed-rank test is
bounded by the number of instances available, and on this benchmark that bound is binding: with
four and three paired instances the smallest attainable two-sided $p$ is $0.125$ and $0.25$, so
at 1{,}000 and 2{,}000 SKUs no result, however unanimous, reaches the conventional 0.05
threshold. Not every cell there is even saturated: the two rows at $p=0.38$ and $p=1.00$, both
at 1{,}000 SKUs, sit above their floor of $0.125$ because the instances are split rather than
unanimous, whereas a cell that does attain its floor is as decisive as its sample permits. At
both sizes we therefore treat the win count as the primary evidence, which is the form in which
the main article states them.

The separations are reported without correction for multiplicity. Of the seven separations
below 0.05 that the table shows, the five at $p \le 0.0020$ survive a Bonferroni, Holm or
Benjamini--Hochberg correction over the eighteen tests reported, whereas the two at
$p = 0.0273$ and $p = 0.049$ do not. A small $p$ indicates that the methods differ, whereas a
large $p$ indicates only that these instances are insufficient to separate them, which is a
weaker statement than equivalence. Effect size and consistency are therefore reported
separately: the Gap column of the main article's benchmark table quantifies the magnitude of a
difference, and the test quantifies how reliably it recurs.

\section{Implementation, hardware and baseline parameters}\label{sup:impl}
This section states in full the computational setup that Section~5.2 of the main article
summarises, together with the rule of each operator of the price-and-complete drive of Section~4.2.3, the incremental swap evaluation, the stage budgets of
the price-and-complete drive, and the parameters of the two literature baselines.

\subsection{Operators of the price-and-complete drive}\label{sup:cgops}
Section~4.2.3 of the main article describes the four operators of its Algorithm~1 in words. This
subsection gives the rule each one applies, with $\Pi(s)$ the products a layout $\Pi$ places at
station $s$, $\sigma_p$ and $\mu$ the duals of the covering and cardinality rows, and $B$ the time
cap.

\textsc{Descend} minimises the penalised visit count
\[
\Psi(\Pi)=\sum_{u \in U} w_u\,\bigl\lvert\{\,s \in S : \Pi(s) \cap u \neq \emptyset\,\}\bigr\rvert
\;+\; \rho \sum_{s \in S} \max\bigl(0,\ \mathrm{load}(s) - T_s\bigr),
\qquad \rho=\sum_{u\in U} w_u,
\]
the visits counted over the distinct multi-item supports plus $\rho$ times the total workload
excess, so that no exchange can lower $\Psi$ by trading feasibility for visits. It exchanges two
products held at different stations, which leaves every station's product count unchanged, and
accepts the first exchange that lowers $\Psi$. Each sweep visits the catalogue in a fresh random
order; a product is paired with every later product when the catalogue holds at most 300, and with
48 partners drawn at random beyond that. The sweep restarts from the improved layout after each
accepted exchange, and the descent stops when a full sweep improves nothing or its budget expires,
$\min(15\,\mathrm{s},\,0.08B)$ on the warm start and $\min(2\,\mathrm{s},\,0.02B)$ on each
completion inside the generation loop. Section~\ref{sup:swapeval} gives the incremental evaluation
that makes one exchange cheap to score.

\textsc{GreedyPrice} returns one bundle per master solve. Its candidates are the products of
strictly positive dual price, ordered in two ways, by decreasing $\sigma_p$ and by decreasing
$\sigma_p$ per unit of workload $L_p/V$. Along each ordering it extends a bundle while the slot and
line limits hold, tracking the exact reduced cost after every addition, and keeps the prefix of
least reduced cost; the better of the two prefixes is returned, and nothing is returned when
neither is negative.

\textsc{ExactPrice} solves the pricing subproblem of Section~4.2.2 of the main article as an
integer programme under a per-call time limit. It returns the bundles found, the solver's dual
bound $\mathrm{rc}_{\mathrm{lb}}$, which remains valid when the call is truncated, and a flag
recording whether the call reached proven optimality.

\textsc{Complete} extends a priced bundle to a full layout. It takes the remaining products in
decreasing $L_p$ and gives each to the station that admits it under both caps and whose products
already touch the greatest weight of supports containing it, ties going to the lowest-indexed
station. A product no station admits under the workload cap goes to the least loaded station,
which can leave the completion infeasible; the gate of Stage~2 then discards it.

The warm start of Stage~1 seeds the pool with the LPT partition and 20 mutated variants of it.
Each variant displaces a random 5 to 10\% of the product-to-station assignments, reinserts the
displaced products in random order at the first station of a random permutation that admits them
under both caps, and drops any product no station admits, so no infeasible column enters the pool.

The pricing model is built once per instance: its constraint matrix does not depend on the duals,
so each call rewrites only the objective coefficients of the $a_p$ variables and, on the
station-indexed master, the two resource limits. The build costs about 21 seconds on a 2,000-SKU
instance; the solves it serves dominate, and the number of pricing iterations reached inside the
cap is what limits the method.

\subsection{Incremental evaluation of the swap descent}\label{sup:swapeval}
The descent of Section~4.2.3 of the main article keeps a support-count matrix
$\mathrm{cnt}[i][u]$, the number of products of station $i$ contained in support $u$. Station $i$
visits $u$ exactly when $\mathrm{cnt}[i][u]>0$. Removing a product frees the supports whose count
falls to zero and inserting it costs those whose count rises from zero. Only supports containing
the two moved products can change, so a swap is scored in time proportional to their support
degrees rather than by re-scanning the orders, and the counts are updated the same way once a swap
is accepted. Beyond about 300 products the candidate partner set is sampled to keep the cost of one
descent step in hand.

\subsection{Stage budgets of the price-and-complete drive}\label{sup:budgets}
The time cap $B$ of Section~4.2.3 of the main article is divided across the four stages of the
drive, and Table~\ref{tab:budgets} gives the two schedules used. The heterogeneous industrial run
is matched to its longer runs: it allots smaller shares to the warm start and the generation loop
than the homogeneous benchmarks do, and it inserts a bound pass that the homogeneous schedule does
not use.

\begin{table}[H]
\centering
\tbl{Stage budgets of the price-and-complete drive, as cumulative fractions of the time cap $B$.}
{\begin{tabular*}{\textwidth}{@{\extracolsep{\fill}}lcc@{}}
\toprule
Stage & Homogeneous benchmarks & Industrial run \\
\midrule
Warm start      & $0.08B$   & $0.05B$   \\
Generation loop & $0.70B$   & $0.62B$   \\
Bound pass      & not used  & $0.74B$   \\
Integer master  & $0.85B$   & $0.82B$   \\
Final polish    & remainder & remainder \\
\bottomrule
\end{tabular*}}
\tabnote{Each entry is the point in the cap at which that stage ends, so the generation loop of the homogeneous schedule runs from $0.08B$ to $0.70B$. The final polish takes whatever the integer master leaves.}
\label{tab:budgets}
\end{table}

\subsection{Baseline operators and parameters}\label{sup:baselines}
Both baselines score a candidate layout with the same surrogate: the station-visit count, plus a
penalty on the total excess over slot capacity or workload cap, plus a dispersion term on the
co-occurrence weight of product pairs placed at different stations. The first term is the real
objective, the second makes an infeasible layout uncompetitive, and the third breaks ties between
layouts the visit count cannot separate. Table~\ref{tab:baselineparams} states the weights of that
surrogate and the operator settings of both procedures.

The genetic algorithm uses binary-tournament parent selection, crossover copying a contiguous
segment of one parent into the other with a repair pass for over-capacity stations, mutation
exchanging two products' stations, and the best layout of each generation carried forward
unchanged. Each initial member is drawn by assigning products to stations uniformly at random
subject to $\zeta_s$, as \citet{pan2015} prescribe.

Simulated annealing with correlated interchange cools geometrically until its stopping temperature
is reached. A move is drawn from the correlated pair list, co-occurring pairs being sorted by
descending co-occurrence and split into equal groups from which a group and then a pair are drawn
uniformly, and the two products exchange stations. A move lowering the surrogate is always
accepted, whereas one raising it by $\Delta$ is accepted with probability $\exp(-\Delta/T)$. The
run starts from a cube-per-order-index (COI) layout, which is correlation-blind like LPT but ranks
products by demand volume per unit of storage occupied rather than by order-line frequency, filling
stations in that order \citep{kim2012}.

\begin{table}[H]
\centering
\tbl{Surrogate weights and operator settings of the two literature baselines.}
{\begin{tabularx}{\textwidth}{@{}Xl@{}}
\toprule
Parameter & Value \\
\midrule
\multicolumn{2}{@{}l}{\emph{Shared surrogate}}\\
Excess penalty on slot capacity and workload cap & $1{,}000$ \\
Dispersion weight on separated correlated pairs  & $0.01$ \\
\midrule
\multicolumn{2}{@{}l}{\emph{Genetic algorithm} \citep{pan2015}}\\
Population size                   & $50$ \\
Maximum generations               & $200$ \\
Crossover probability             & $0.8$ \\
Crossover segment cap             & $0.25\,|P|$ \\
Mutation probability              & $0.1$ \\
\midrule
\multicolumn{2}{@{}l}{\emph{Simulated annealing with correlated interchange} \citep{kim2012}}\\
Initial temperature               & $800$ \\
Geometric cooling ratio per level & $0.95$ \\
Stopping temperature              & $1$ \\
Moves proposed per level          & $\min(2|P|,\,1000)$ \\
Correlated-pair groups            & $10$ \\
\bottomrule
\end{tabularx}}
\tabnote{Both procedures also stop at the per-size wall-clock cap of Section~5.2 of the main article, whichever comes first. The crossover segment cap of $0.25\,|P|$ is a quarter of the catalogue. Seeds and solver versions are given in Section~\ref{sup:hardware}.}
\label{tab:baselineparams}
\end{table}

\subsection{Hardware, solver versions and seeds}\label{sup:hardware}
Every run reported in the main article was produced on the same remote machine, a 16-core Windows
Server 2019 host with 128\,GB of memory, under the same Python 3.10 environment, and every method
received the same core and thread allocation, four threads, so no method was given more compute
than another under the shared wall-clock cap. The column generation runs on CPLEX 22.1.1 and the
two MILP representations on Hexaly 13.0. One run per instance is reported in every cell. The two
baselines are our own Python implementations of the published operators, so a residual
implementation-efficiency difference against the two commercial engines cannot be excluded, and we
record how far each baseline got rather than assume it exhausted its search.

Both metaheuristic baselines draw their search decisions from a NumPy generator seeded at 42, so
SA-C is fully determined by the instance and the cap. The genetic algorithm's capacity repair draws
instead on the global random stream, seeded at 20240612 by the experiment harness, and that stream
is the only source of run-to-run variation in the reported GA figures. The column-generation drive
is also stochastic, since Stage 1 mutates the warm start and the descent samples candidate partners
beyond about 300 products, and it draws on the same harness stream. We report one run per cell for
it as well, so the 2,000-SKU result rests on three single draws and a replicated cell would
strengthen it.

\section{Proofs of the bounding claims}\label{sup:proofs}

\subsection{Binary MILP linear relaxation floor}\label{sup:bounds_milp}
Take any point feasible for the linear relaxation. For an order $o \in O$ and a product
$p \in P_o$, the linking row $x_{ps}\le z_{os}$ forces $z_{os}\ge x_{ps}$ for every station, and
summing over the stations together with the covering row $\sum_{s\in S} x_{ps}\ge 1$ gives
$\sum_{s\in S} z_{os}\ge\sum_{s\in S} x_{ps}\ge 1$. The objective
therefore cannot fall below $|O|$ at any feasible point, fractional or integral. That floor is
attained: setting $x_{ps}=1/|S|$ meets the capacity rows whenever $|P|/|S|\le\zeta_s$ and the
workload rows under the operational slack of Section~5.1 of the main article, so $z_{os}=1/|S|$ is
feasible and the objective equals $|O|$ exactly. On this instance family the relaxation is
therefore worth precisely the statement that every order visits at least one station, independent
of the co-occurrence structure that makes one instance harder than another, so a branch-and-bound
search starts from that value with nothing instance-specific to prune with.

\subsection{Column generation Farley-type dual bound}\label{sup:bounds_cg}
The bound does not rest on the price-and-complete drive. It follows from linear-programming duality
applied to the restricted master together with a valid dual bound from the pricing solve, and would
hold whatever heuristic produced the layouts.

Section~4.2.1 of the main article aggregates the station-indexed master over interchangeable
stations. Written over the current pool $Q'$ of feasible bundles, with $c_q$ the weighted-support
cost of bundle $q$, that restricted master is
\begin{align}
\min \quad & \sum_{q \in Q'} c_q\, y_q \notag\\
\text{s.t.} \quad
& \sum_{q \ni p} y_q \ge 1 && \forall p \in P \quad (\sigma_p \ge 0) \label{sup:eq:cover}\\
& \sum_{q \in Q'} y_q \le |S| && (\mu \le 0) \label{sup:eq:card}\\
& 0 \le y_q \le 1 && \forall q \in Q'. \notag
\end{align}
The upper bounds $y_q\le 1$ may be dropped, since reducing any $y_q$ above one to one preserves
coverage, relaxes the cardinality row and cannot increase the objective, so an optimal dual
solution with zero multipliers on them exists. Strong duality gives
$z_{\mathrm{RMP}}=\sum_p\sigma_p+|S|\mu$. The pair $(\sigma,\mu)$ need not be dual-feasible for the
full master, because a bundle outside the pool may price out negatively. Let $\mathrm{rc}_{\min}$
be the least reduced cost over all feasible bundles and suppose it is negative. Replacing $\mu$ by
$\mu'=\mu+\mathrm{rc}_{\min}$ restores feasibility, since every bundle then satisfies
$c_q-\sigma(q)-\mu'=\mathrm{rc}_q-\mathrm{rc}_{\min}\ge 0$, and $\mu'$ remains non-positive. Weak
duality at this repaired point gives
\begin{equation*}
z_{\mathrm{LP}}\ge\sum_p\sigma_p+|S|\mu'=z_{\mathrm{RMP}}+|S|\cdot\mathrm{rc}_{\min}.
\end{equation*}
Substituting the pricing solve's own dual bound $\mathrm{rc}_{\mathrm{lb}}$, which never exceeds
$\mathrm{rc}_{\min}$, only weakens the inequality, and taking $\min(0,\cdot)$ leaves the bound at
$z_{\mathrm{RMP}}$ once the master has priced out. Because the covering rows relax the partition
and the partition model reproduces the visit count exactly, the chain
$\mathrm{LB}\le z_{\mathrm{LP}}\le z_{\mathrm{IP}}\le$ the optimal visit count holds.

Stated in the form the main article reports, the bound is
\[
\mathrm{LB} = z_{\mathrm{RMP}} + |S|\cdot \min\bigl(0,\ \mathrm{rc}_{\mathrm{lb}}\bigr)
\]
on the aggregated master and
\[
\mathrm{LB} = z_{\mathrm{RMP}} + \sum_{s \in S} \min\bigl(0,\ \mathrm{rc}^{\mathrm{lb}}_s\bigr)
\]
on the station-indexed one, the same argument being applied to each convexity row in turn. Within
the caps of the benchmark the bound is informative only at 50 SKUs, where the mean bound of 3,312
exceeds the trivial floor of one visit per order by 63\% yet sits well below the mean of 7,590
visits found; from 500 SKUs upward it is vacuous, for the reason Section~\ref{sup:cgrecord}
records.

Three conditions make the argument valid and the implementation meets all three. Pricing is solved
as an exact integer programme, so its dual bound holds even when a call is truncated. The
station-indexed repair is applied one convexity row at a time at a common dual vector, which the
run enforces after the generation loop. The bound covers the multi-item supports, to which the
singleton constant of Section~4.2.1 of the main article is added. Degeneracy of the covering master
slows how quickly the floor rises as columns arrive but never invalidates it, and we declare
convergence only when an exact solve at the true duals proves that no bundle of negative reduced
cost remains.

\section{The Company~A instance}\label{sup:instance}
Section~6 of the main article gives the decision pool and the workload basis in prose. This section
records the dataset behind them, the residual resources the optimisation actually receives, and the
two parameters the site fixes rather than we do.

\begin{table}[H]
\centering
\tbl{The Company~A instance as it enters the solvers.}
{\small\begin{tabular*}{\textwidth}{@{\extracolsep{\fill}}lr@{}}
\toprule
Quantity & Value \\
\midrule
Planning horizon & 13 weeks, 1 Sept.\ to 30 Nov.\ 2021 \\
Orders & 284{,}862 \\
Order lines over the horizon & 1{,}486{,}608 \\
Products (SKUs) & 21{,}874 \\
\quad free to relocate & 15{,}975 \\
\quad frozen in place & 5{,}899 \\
Stations, of which evaluated & 26, of which 24 \\
Movable products held per station & 8 to 2{,}575 \\
Processing rate $V_s$, lines per unit time & 6.2 to 10{,}400 \\
Legacy station load $\sum_{p} L_p/V_s$ & 0.8 to 9{,}923 \\
Workload budget $T_s$ & 110\% of the legacy load \\
\bottomrule
\end{tabular*}}
\tabnote{The two excluded stations hold 1 and 2 locations and carry 5 and 11 order lines between them. Processing rates span three orders of magnitude because the stations use different picking technologies, which is why workload is normalised by $V_s$ and raw line counts are not comparable across stations.}
\label{tab:instance}
\end{table}

Two parameters are the site's operating rules rather than choices made here. A product is frozen
when it appears in at most five orders \emph{and} sits at one of the three static-shelving
stations, the five-order threshold being the site's own; and the workload budget $T_s$ is $110\%$
of a station's complete legacy load, the $10\%$ being the tolerance the site operates to. Neither
is tuned here, which is why no sensitivity is reported over them.

Frozen products keep their slots and their workload share, so the optimisation runs on residual
resources rather than full ones. Writing $P^{\mathrm{fix}}_s$ for the products frozen at station
$s$,
\[
\bar\zeta_s=\zeta_s-\lvert P^{\mathrm{fix}}_s\rvert,
\qquad
\bar T_s = T_s - \sum_{p \in P^{\mathrm{fix}}_s} L_p/V_s,
\]
and the solver receives $\bar\zeta_s$ and $\bar T_s$ in place of $\zeta_s$ and $T_s$. Evaluation
reverses this: every station is charged with its complete load, frozen and movable products
together, over the complete stream of $284{,}862$ orders rather than the $83{,}183$ that contain a
movable product. The reductions of the main article's industrial table are therefore measured
against the layout the site operates, on the whole catalogue, and are not inflated by the reduced
decision pool.

\section{The industrial community-bound search}\label{sup:betasearch}
Section~6 of the main article selects the community bound $\beta$ for Company~A, whose stations
hold between 8 and 2{,}575 movable products, on a seven-point grid rather than from the rule of
Section~4.1. This section records the full search behind that choice.

Applying the rule $\beta=\max(5,\lfloor 0.4\,\zeta+\tfrac12\rfloor)$ to the mean station capacity
returns $\beta=266$. That setting is feasible and removes $5.66\%$ of the visits the site incurs
today, so the rule does yield a deployable layout here. It is simply not the best value reachable,
and nothing in the rule reveals as much.

Table~\ref{tab:betasearch} reports the fifteen settings evaluated. Feasibility is not an interval
in $\beta$, and the shape it does take is the one the split rule predicts. A community larger than
a station has to be split at placement, and the split rule then places products one at a time under
the same admissibility test, so the blocks it produces are sized by what the stations can take
rather than by $\beta$. Both ends of the range therefore respect the small stations: the low end
because its communities already fit them, the high end because almost every community is split
before placement. The middle is the worst place to sit, with communities too large for the small
stations and too small to force the split path throughout. The four infeasible settings, from 33 to
133, are exactly that middle.

The protocol a site would run is the grid anchored on the mean capacity,
$\{33,67,133,200,266,333,399\}$, keeping the feasible layout with fewest visits. Five of the seven
are feasible and $\beta=333$ is the best of them, which gives the heuristic row of the main
article's industrial table and is why that row's time entry charges the whole grid, 790 seconds.
The other eight settings were exploratory and are excluded from that cost; charging all fifteen
would raise the heuristic to $1{,}480$ seconds, still an order of magnitude below the next fastest
method. One exploratory setting is better, $\beta=300$ at $994{,}596$ visits against $996{,}448$, a
further $0.19\%$, but a grid specified in advance does not contain it.

What the grid buys is not a finely tuned bound but the avoidance of the infeasible band. From
$\beta=200$ upward every setting evaluated is feasible, and those nine span $5.34\%$ to $6.39\%$,
a band of one percentage point: once the bound is large enough to force the split path the visit
count is insensitive to it. Nothing therefore rests on the particular value the grid returns, and
we do not read $\beta=333$ as a calibrated quantity. The insensitivity is not total, because the
settings differ more in how they leave the workload than in what they save: $\beta=450$ reaches
$6.31\%$ at a utilisation dispersion of $21.6$ against the selected setting's $27.4$. We keep the
selection rule as stated rather than adding dispersion to it, since the rule was fixed before these
runs and choosing on a criterion the outcome suggested would not be a protocol a site could follow
in advance. On a warehouse of interchangeable stations none of this arises, the rule's $\beta$
being feasible on all 29 benchmark instances.

\begin{table}[H]
\centering
\tbl{The industrial community-bound search on the Company~A instance.}
{\small\begin{tabular*}{\textwidth}{@{\extracolsep{\fill}}rrrrl@{}}
\toprule
$\beta$ & Visits & Reduction (\%) & Time (s) & Status \\
\midrule
15  & 1{,}037{,}244 & 2.38 & 50  & feasible \\
33  & 1{,}010{,}119 & 4.93 & 97  & infeasible, grid \\
47  & 1{,}000{,}007 & 5.88 & 79  & infeasible \\
67  & 996{,}734     & 6.19 & 105 & feasible, grid \\
100 & 1{,}001{,}182 & 5.77 & 74  & infeasible \\
133 & 1{,}007{,}309 & 5.20 & 123 & infeasible, grid \\
200 & 1{,}001{,}428 & 5.75 & 89  & feasible, grid \\
250 & 1{,}005{,}780 & 5.34 & 78  & feasible \\
266 & 1{,}002{,}379 & 5.66 & 87  & feasible, grid, rule \\
300 & 994{,}596     & 6.39 & 70  & feasible \\
333 & 996{,}448     & 6.22 & 94  & feasible, grid, selected \\
350 & 1{,}000{,}162 & 5.87 & 119 & feasible \\
365 & 1{,}001{,}206 & 5.77 & 117 & feasible \\
399 & 1{,}000{,}199 & 5.86 & 194 & feasible, grid \\
450 & 995{,}426     & 6.31 & 104 & feasible \\
\bottomrule
\end{tabular*}}
\tabnote{Reduction is against the legacy layout's $1{,}062{,}507$ visits. A setting is feasible when no station exceeds its own workload budget, $110\%$ of its complete legacy load; the check is per station, so a low visit count does not imply feasibility. ``grid'' marks the seven settings of the protocol a site would run, ``rule'' the value returned by $\beta=\max(5,\lfloor 0.4\,\zeta+\tfrac12\rfloor)$ on the mean station capacity, and ``selected'' the feasible grid setting with fewest visits, which gives the heuristic row of the main article's industrial table.}
\label{tab:betasearch}
\end{table}

\section{Per-instance column-generation convergence record}\label{sup:cgrecord}
The released archive named in the main article's data availability statement carries, for every
column-generation run, the convergence status, the restricted-master value $z_{\mathrm{RMP}}$, the
pricing dual bound $\mathrm{rc}_{\mathrm{lb}}$ returned by the last pricing call, and the deadline
the generation loop reached. Those records support the reading of Section~4.2.4 of the main
article: no run proved that no improving bundle remained, the generation loop reached its $0.70B$
allocation on all 29 instances at all four sizes, and the last dual bound was strictly negative
everywhere, from $-696$ to $-1{,}497$ at 50 SKUs and widening to $-2.3\times10^{5}$ at 2,000 SKUs.
The quantity that pushes the reported bound down is therefore the correction
$|S|\cdot\mathrm{rc}_{\mathrm{lb}}$ rather than any weakness of the relaxation, and the binding
constraint is the number of iterations reached inside the cap.

\bibliographystyle{apacite}
\bibliography{IJSSOL_CSLAP_v1}

\end{document}
